\chapter{The Einstein Equations}
\label{chap:the-einstein-equations}

% Closed question: What records the failure of transport to commute?

\begin{chapteressay}

A car turns through a tight curve. The inner wheel travels a shorter path than
the outer wheel. If the speedometer reads only one wheel, the odometer can be
slightly wrong during the turn. The car has not stopped measuring. It has
entered geometry.

This is the chapter's everyday doorway into curvature. On a straight road, it
is easy to imagine one distance count standing for the whole vehicle. In a
turn, different parts of the car trace different arcs. Transport is no longer
transparent. The mismatch is not an error added after the fact. It is the
physical signature of the path.

General relativity makes this severe. Carry a gyroscope around Earth and its
axis does not return exactly as flat space would predict. Let two long
interferometer arms compare their lengths while a gravitational wave passes,
and the differential strain records spacetime curvature moving through the
detector. Watch stars orbit far from a galactic center, and their speeds do not
fall as the visible mass budget predicts. Something in the curvature budget is
missing.

The same structure appears outside gravity. In a turbulent river, transporting
smooth velocity data forward can fail because gradients feed smaller gradients.
In neutrino oscillation, a flavor state transported through space does not
remain an electron-flavor record. The measured state changes according to a
mass-basis structure that does not commute with the flavor basis.

The code names this chapter with \lean{SpaceTimePath}, \lean{CalculusProcess},
\lean{UNIVERSAL}, \lean{YarnTheory}, and \lean{HeartbeatProcess}. A path can be
Einstein, white hole, blackhole, or geodesic. A yarn can be a Stokes record, a
fiber comparison, or a fabric of several paths. The names are satire, but the
shape is exact: transport has grammar, comparison has order, and accumulated
strain becomes a record.

The closed question is:

\begin{center}
\emph{What records the failure of transport to commute?}
\end{center}

The answer is curvature, read as strain in the measurement ledger. If carrying
structure along one path and then another gives a different result from the
opposite order, the mismatch is not noise. It is what the field is measuring.

Chapter 9 therefore turns calibration into dynamics. Chapter 8 asked how
observers compare counts. Chapter 9 asks what the comparison records when
transport itself has path-dependent residue.

\end{chapteressay}

% Phenomenon arcs use: 1/3 setup -> phenomenon box -> 2/3 structural interpretation.

\section{Strain as Transport Failure}
\label{se:strain-as-transport-failure}

Gravity Probe B was launched on April 20, 2004, into a $642\,\mathrm{km}$
polar orbit around Earth. The spacecraft carried four gyroscopes machined
from fused quartz to the most precise sphericity ever achieved: deviations
from a perfect sphere of less than $10\,\mathrm{nm}$ over a $38\,\mathrm{mm}$
diameter ball. Each gyroscope was coated with a thin layer of niobium and
cooled to $1{.}8\,\mathrm{K}$ in a helium dewar, making it
superconducting. The spinning superconductors produced magnetic moments
aligned with their rotation axes; superconducting quantum interference
devices (SQUIDs) measured the magnetic moments and hence the orientations
of the spin axes to within $10^{-5}$ arcseconds. The spacecraft also
carried a telescope locked onto the guide star IM Pegasi (HR 8703), so
that the gyroscope orientations could be measured relative to a fixed
distant reference.

The mission's goal was to measure two relativistic effects predicted by
general relativity: the \emph{geodetic effect} (de Sitter precession),
which arises from the curvature of spacetime caused by the Earth's mass;
and the \emph{frame-dragging effect} (Lense-Thirring precession), which
arises from the Earth's rotation dragging the surrounding spacetime
along with it. The predicted magnitudes were
\[
\Omega_{\text{geodetic}} = 6{,}606 \,\mathrm{mas/yr}, \qquad
\Omega_{\text{frame-drag}} = 37{.}2 \,\mathrm{mas/yr},
\]
where ``mas'' denotes milli-arcseconds. The geodetic effect points in the
orbital plane; the frame-dragging effect points perpendicular to it.

Mission operations ran from August 2004 through August 2005, with data
analysis continuing for another six years due to unexpected systematic
effects from the gyroscopes' patch effects (electrostatic anomalies on the
sphere surfaces). The final published result in 2011 reported the geodetic
precession as $6{,}601{.}8 \pm 18{.}3\,\mathrm{mas/yr}$ and the
frame-dragging precession as $37{.}2 \pm 7{.}2\,\mathrm{mas/yr}$. Both
values agreed with the general relativistic predictions within the
experimental uncertainty.

The structural significance is what the gyroscope is doing. It is not
measuring gravity as a force pulling on a pointer; gyroscopes do not
respond to forces in the relevant sense. It is measuring how an orientation
changes when transported along a worldline through curved spacetime.

In flat spacetime, ideal parallel transport of a vector around a closed
loop returns the vector unchanged. The orientation of an isolated gyroscope
in inertial motion through flat spacetime would not precess relative to
the distant stars; it would always point the same way. In curved spacetime,
parallel transport around a closed loop does not return the vector
unchanged: the vector rotates by an amount determined by the curvature
enclosed by the loop. A gyroscope in orbit around the Earth traces a
closed loop (or, more precisely, a tightly wound helix; over many orbits
it samples a closed loop's worth of curvature each year). The gyroscope's
orientation, observed against a distant fixed reference, precesses by an
amount that records the integrated curvature.

This is the chapter's central claim. Curvature is not an invisible
substance suffusing space. Curvature is the structured record of
transport's failure to commute or to return unchanged. Gravity Probe B
measured curvature by measuring the rotation accumulated during transport.
The measurement was direct: the gyroscope orientation was a physical
quantity recorded continuously throughout the mission; the precession
was a real, measured rotation; the rotation rate matched the
predictions of general relativity.

\begin{phenom}{Gravity Probe B Effect}
\label{ph:gravity-probe-b-effect}

\PhStatement Curvature can be measured as the mismatch left after transporting
an orientation around a physical path.

\PhOrigin Gravity Probe B tested general relativity by tracking gyroscope
precession in a polar orbit around Earth.

\PhObservation The gyroscope axes accumulated geodetic and frame-dragging
precession relative to a distant guide star, matching the relativistic
transport prediction within experimental uncertainty.

\PhConstraint The measurement may not treat orientation as globally fixed when
the spacetime path itself changes the transported frame.

\PhConsequence Curvature is not an invisible substance. It is the recorded
failure of transport to return structure unchanged.

\end{phenom}

This is the operational heart of the Einstein equations. Geometry is what
transport does. The connection tells how to carry a vector locally.
Curvature is the residue of carrying it around a loop. The Einstein
equations then balance that curvature record against stress-energy.

Gravity Probe B is the cleanest direct test of this picture. The
gyroscopes were nearly perfect: any precession not predicted by
relativity would have been visible at the milli-arcsecond level. The
guide star provided a fixed external reference: the gyroscope orientations
were measured against IM Pegasi, whose own motion relative to the
extragalactic frame was independently determined by VLBI observations
to better than $0{.}1\,\mathrm{mas/yr}$. The redundancy of four
gyroscopes allowed systematic effects to be calibrated against the
expected pattern. The result is a measurement of curvature as a
physical quantity --- not an inferred theoretical construct, but a
recorded rotation.

The chapter's structural framework is precise here. The gyroscope's
spin axis is a record that is transported along the spacecraft's
worldline. In a flat background, the transport would preserve the
axis direction. In the Schwarzschild-plus-Lense-Thirring background
of a rotating massive Earth, the transport precesses the axis. The
precession is the curvature: it is the algebraic record of how much
the geometry deviated from flatness along the path. The Einstein
field equations relate this curvature to the stress-energy content
of spacetime (the Earth's mass-energy and angular momentum); Gravity
Probe B measured the curvature side of the equation directly.

The Lean formalization in \lean{HeartbeatProcess} captures the same
idea at the algebraic level. A \lean{HeartbeatProcess} accumulates
elaboration cost along a proof path; two paths to the same
conclusion may accumulate different costs; the difference is the
strain of the path. For Gravity Probe B, the analog is: the gyroscope
records the integrated curvature along its worldline; two worldlines
with the same endpoints (e.g., one full orbit, vs. staying at rest
relative to the distant stars) accumulate different precession; the
difference is the curvature enclosed. The strain is physical in the
gyroscope case and algebraic in the proof case, but the structure
is the same: transport failure is what curvature looks like as a
recorded quantity.

\section{The Navier--Stokes Third Parameter}
\label{se:the-navier-stokes-third-parameter}

A river can look smooth at one scale and violent at another. Velocity changes
from point to point. Gradients feed gradients. Eddies break into smaller
eddies. A model that transports the velocity field forward must account for
whether the smooth record remains smooth.

The Navier-Stokes equations are not Einstein's equations, but they belong in
this chapter because they ask a parallel measurement question: can the
transport of a field remain globally controlled, or can finite-time strain
make the record fail?

Write the incompressible equation schematically as

\[
\partial_t v + (v \cdot \nabla)v = -\nabla p + \nu \Delta v.
\label{eq:chapter9-navier-stokes}
\]

The nonlinear transport term carries velocity by velocity. The viscous term
smooths. The conflict between them is the source of the problem.

\begin{phenom}{Navier-Stokes Effect}
\label{ph:navier-stokes-effect}

\PhStatement A transported field can generate strain faster than smoothing can
control it, making global continuation itself the measurement question.

\PhOrigin The Navier-Stokes existence and smoothness problem asks whether
smooth three-dimensional incompressible flow can develop singular behavior in
finite time.

\PhObservation Turbulent flows transfer structure across scales, so small
errors or gradients can become dynamically significant rather than remaining
local bookkeeping details.

\PhConstraint The model may not assume a smooth global record merely because
the initial data were smooth.

\PhConsequence Field transport must carry its own strain budget; continuation
is admissible only while that budget remains controlled.

\end{phenom}

This section is deliberately adjacent rather than identical to relativity. The
point is not that fluids are spacetime. The point is that transport can fail to
commute as an accounting matter. Carry local data forward, compare it at a
later scale, and the mismatch becomes the phenomenon.

\section{Flat Rotation Curves}
\label{se:flat-rotation-curves}

In a galaxy whose mass is mostly concentrated near the center, orbital speeds
should fall with radius, roughly like $1/\sqrt r$ outside the visible disk. But
many spiral galaxies show flat rotation curves: stars far from the center move
about as fast as stars nearer in.

The record does not balance if only visible matter is counted.

\begin{phenom}{Flat Rotation Curve Effect}
\label{ph:flat-rotation-curve-effect}

\PhStatement A curvature or mass budget can fail operationally when transported
orbital records do not match the visible source count.

\PhOrigin Measurements of spiral galaxy rotation curves showed orbital speeds
remaining approximately flat far beyond the luminous disk.

\PhObservation The observed velocities imply more gravitational mass than is
seen in stars and gas, or else a change in the gravitational accounting rule.

\PhConstraint The measurement may not close the curvature ledger using only
visible matter when the orbital records demand additional source structure.

\PhConsequence Dark matter, modified gravity, or any rival explanation must pay
the same bill: the transported orbital record has a missing source.

\end{phenom}

This example gives the Einstein-equation idea its public shape. The left side
of the equation is geometry. The right side is stress-energy. If the geometry
inferred from motion does not match the visible stress-energy, the ledger is
open. Something has not been counted.

\section{Gravitational Waves from LIGO}
\label{se:gravitational-waves-from-ligo}

The Laser Interferometer Gravitational-Wave Observatory (LIGO) consists of
two facilities, one in Hanford, Washington, and one in Livingston, Louisiana,
separated by about $3{,}000\,\mathrm{km}$. Each facility is an L-shaped
Michelson interferometer with $4\,\mathrm{km}$ arms. A laser beam is split at
the corner of the L; each half travels down one arm, reflects off a hanging
mirror at the far end, and returns to the corner; the two returning beams
are recombined and the interference pattern is recorded by a photodiode.

LIGO measures strain:
\[
h = \frac{\Delta L}{L}.
\label{eq:chapter9-ligo-strain}
\]
A passing gravitational wave stretches one arm while compressing the other,
then reverses. The strain $h$ is the fractional change in arm length. For
$L = 4\,\mathrm{km}$ and a typical astrophysical source, the strain at
Earth is on the order of $10^{-21}$, meaning the arm length changes by
$\Delta L = 10^{-21} \times 4 \times 10^3 = 4 \times 10^{-18}\,\mathrm{m}$,
or about one-thousandth the diameter of a proton.

The signal is fantastically small. To detect it, the interferometer must
suppress noise sources orders of magnitude larger. Seismic noise from
ground vibrations is suppressed by a four-stage pendulum suspension
system that decouples the mirrors from the ground at frequencies above
$10\,\mathrm{Hz}$. Thermal noise in the mirror coatings is suppressed by
operating the mirrors as cold as the cryogenic infrastructure allows
(LIGO operates at room temperature; KAGRA, a Japanese sister
instrument, operates cryogenically). Shot noise in the laser is
suppressed by using a high-power laser (about $750\,\mathrm{kW}$
circulating in each arm cavity) and by injecting squeezed vacuum states
into the dark port. Quantum radiation pressure noise at low frequencies
is balanced against shot noise at high frequencies through the choice
of laser power.

The result is a strain sensitivity that reaches about $10^{-23}$ per
$\sqrt{\mathrm{Hz}}$ in the most sensitive band around
$100\,\mathrm{Hz}$. Integrated over the duration of a typical
astrophysical signal (a fraction of a second), the detectable strain is
in the $10^{-22}$ to $10^{-21}$ range.

On September 14, 2015, at 09:50:45 UTC, the LIGO detectors at Hanford
and Livingston both recorded a signal consistent with the merger of two
black holes. The signal lasted about $0{.}2\,\mathrm{s}$ and swept
upward in frequency from $35\,\mathrm{Hz}$ to $250\,\mathrm{Hz}$ --- a
characteristic ``chirp'' as the orbital frequency increased while the
two black holes spiraled together. The peak strain was about $10^{-21}$.
The peak amplitude corresponds to the moment of merger. The signal then
faded as the merged black hole rang down to its final stationary
configuration. This was GW150914, the first direct detection of
gravitational waves.

Analysis of the waveform yielded the source parameters: two black holes
of masses approximately $36$ and $29$ solar masses, merging at a distance
of about $1{.}3$ billion light-years (redshift $z \approx 0{.}09$). The
final black hole had a mass of $62$ solar masses, meaning that
approximately $3$ solar masses' worth of energy ($3 M_\odot c^2 \approx
5{.}4 \times 10^{47}\,\mathrm{J}$) was radiated away as gravitational
waves during the merger --- in less than a second. The peak luminosity
exceeded $3{.}6 \times 10^{49}\,\mathrm{W}$, more than the total
luminosity of all the stars in the observable universe.

The detection relied on coincidence between the two LIGO sites. The
Hanford signal arrived $7\,\mathrm{ms}$ before the Livingston signal,
consistent with a gravitational wave traveling at the speed of light
from a source in the southern sky. The waveforms at the two sites,
corrected for the different antenna response patterns, matched to within
the noise. Independent analyses by separate teams using separate
software pipelines confirmed the detection. The statistical significance
was greater than $5{.}1\sigma$.

\begin{phenom}{LIGO Strain Effect}
\label{ph:ligo-strain-effect}

\PhStatement A gravitational wave is measured as differential strain in a
calibrated comparison of spacetime paths.

\PhOrigin Laser interferometer gravitational-wave detectors were built to
compare long optical paths for the passing strain predicted by general
relativity.

\PhObservation GW150914 matched the waveform expected from a binary black-hole
merger, with correlated signals in separated detectors.

\PhConstraint The apparatus may not claim a wave from a single local wiggle; it
must carry differential arm comparison, noise modeling, and cross-detector
coincidence.

\PhConsequence Curvature can be a traveling record. The failure of distances to
remain mutually fixed is the signal itself.

\end{phenom}

LIGO also folds Chapter 5 back into Chapter 9. Interferometry counted phase
distinctions as distance. Here the same technique records spacetime strain.
The detector is not watching a wave pass by like a ripple on a pond. It is
asking whether two calibrated paths remain comparable in the same way.

The structural point is that gravitational waves are recorded curvature
travelling at the speed of light. The Einstein field equations admit wave
solutions in vacuum: $\Box h_{\mu\nu} = 0$, where $h_{\mu\nu}$ is the
linearized metric perturbation. These solutions propagate at $c$ and carry
energy and momentum. When such a wave passes through the LIGO detector,
the strain $h$ at the detector is the metric perturbation transversely
projected onto the arm directions. The chapter's claim is that this strain
is the same kind of object as the gyroscope precession of Gravity Probe B:
both are records of transport failure, both are physical readings of
curvature, both confirm that curvature is a measurable quantity rather
than an abstract theoretical construct.

GW150914 was the first direct detection, but by 2024 LIGO and its
collaborators (Virgo in Italy, KAGRA in Japan) had detected over $100$
gravitational wave events. Most are binary black hole mergers; several
are binary neutron star mergers; a few are mixed black-hole/neutron-star
mergers. The catalog of detections has become large enough to support
population studies: the distribution of black hole masses, the rate of
mergers per unit cosmological volume, the spins and inclinations of the
merging systems. The 2017 detection of GW170817, a binary neutron star
merger, was accompanied by electromagnetic counterparts (a short gamma-ray
burst, a kilonova in optical and infrared, X-ray and radio afterglows)
that pinpointed the source's host galaxy and established gravitational
wave astronomy as a multi-messenger discipline.

The chapter's framework treats the LIGO record as the canonical example
of curvature as a transported quantity. The strain $h(t)$ recorded over
$0{.}2\,\mathrm{s}$ during GW150914 is a time series of metric
perturbation amplitudes; the time series is the curvature wave's
fingerprint passing through the detector; the fingerprint matches the
prediction of general relativity for the specific source. The record is
admissible because the coincidence between sites, the matched-filter
analysis against theoretical waveforms, and the noise model together
support the inference that the signal is a real astrophysical event
rather than an instrumental artifact.

\section{Neutrino Oscillations}
\label{se:neutrino-oscillations}

Neutrinos are produced and detected in flavor states: electron, muon, and tau.
But propagation is governed by mass states. The flavor basis and mass basis are
not the same. As a neutrino travels, the phase differences among mass
components change, so the flavor measured later can differ from the flavor
produced.

\begin{phenom}{Neutrino Oscillation Effect}
\label{ph:neutrino-oscillation-effect}

\PhStatement A carried quantum state can change its measured label when the
transport basis and the observation basis do not commute.

\PhOrigin Solar, atmospheric, reactor, and accelerator neutrino experiments
showed that neutrinos change flavor during propagation, implying nonzero mass
differences.

\PhObservation Electron neutrinos produced in the Sun arrive at Earth as a
mixture of flavors rather than as a conserved electron-neutrino count.

\PhConstraint The measurement may not assume that a label assigned at
production remains invariant under transport when the physical basis differs.

\PhConsequence Noncommuting transport and observation bases create a real
record of change, not a mere naming convention.

\end{phenom}

This closes the chapter by widening curvature into commutation failure. In
gravity, transport around a loop changes orientation. In fluid dynamics,
transport can amplify strain. In galaxies, transported orbital records expose a
source mismatch. In LIGO, differential transport is the wave. In neutrinos, the
failure is between bases.

The common structure is \lean{YarnTheory}: paths, fibers, fabric, and the order
relations that say which carried records can be compared.

\begin{bridge}

The chapter's question was what records the failure of transport to commute.

The answer is strain. Gravity Probe B measures orientation mismatch. Navier-
Stokes asks whether transported smoothness can survive nonlinear strain. Flat
rotation curves expose a missing source in the curvature budget. LIGO records
differential spacetime strain. Neutrino oscillations show a label changing
because transport and observation bases do not commute.

The next chapter asks what remains invariant when labels can change. If
transport can twist records, then physics needs a rule for what survives
admissible relabeling.

\end{bridge}

\begin{coda}{A Bridge Deck with Expansion Joints}

The bridge crosses a river in steel segments. In the morning shade, the deck is
shorter. By afternoon, sunlight warms one span more than another. The steel
expands. The expansion joints open and close.

The joints are not decorative gaps. They are instruments. Each gap records the
failure of the deck to behave as one rigid length under changing heat. If the
engineer pretended every segment expanded identically, the bolts would complain
first, then the welds, then the traffic report.

A truck crossing the bridge transports load from segment to segment. On a cool,
even day, the transfer is quiet. On a hot day with uneven expansion, the load
meets strain at the joints. The path across the bridge still exists, but the
transport is no longer trivial. The joint gap is the visible residue.

This performs the chapter's structure. The deck segment is a local frame. The
joint is a comparison. Thermal expansion is curvature-like strain. The truck's
route is \lean{SpaceTimePath}. The stress analysis is \lean{YarnTheory}. The
sensors watching the gap through the day are \lean{HeartbeatProcess}. The
question is not whether the bridge exists. The question is whether transported
load returns the same account after the path has crossed all those strained
interfaces.

If a joint sticks, the record changes. Expansion that should have been absorbed
becomes force. Force becomes a crack. The crack becomes a maintenance entry.
No one needs a metaphysical theory of bridges to see the rule: transport
failure writes itself somewhere.

The same is true in the physical chapters. A gyroscope writes curvature into
precession. A galaxy writes missing mass into rotation speed. A gravitational
wave writes strain into interferometer arms. A neutrino writes basis mismatch
into flavor statistics.

The bridge deck teaches the quiet version. Curvature is not always a spectacle.
Sometimes it is the extra millimeter in the joint at noon.

\end{coda}
